\documentclass[11pt]{article}
\usepackage[T1]{fontenc}
\usepackage{amsmath,amssymb}
\usepackage{geometry}
\usepackage[hidelinks]{hyperref}
\geometry{margin=1in}

\title{Literature-Implied Answer: Strict \(H^\infty\)-Calculus Angle}
\author{Run \texttt{fa\_banach\_001}, agent \texttt{agent\_lane\_09}}
\date{2026-06-14}

\begin{document}
\maketitle

\section*{Source Question}

The source paper is N. J. Kalton and L. Weis, \emph{The \(H^\infty\)-calculus
and sums of closed operators}, arXiv:math/0010155.  In Section 4, after defining
the sectorial angle \(\omega(A)\) and the bounded \(H^\infty\)-calculus angle
\(\omega_H(A)\), the authors write:
\[
  \text{Is } \omega_H(A)=\omega(A) \text{ whenever } A
  \text{ is } H^\infty\text{-sectorial?}
\]
The sentence appears on page 16 of the source PDF.

\section*{Supporting Theorem}

The supporting paper is N. J. Kalton, E. Lorist, and L. Weis, \emph{Euclidean
structures and operator theory in Banach spaces}, arXiv:1912.09347.  Section
6.4 records the known counterexample lineage: Cowling--Doust--McIntosh--Yagi
first constructed a bounded \(H^\infty\)-calculus example without dense range
with \(\omega(A)<\omega_{H^\infty}(A)\), and Kalton later constructed a
sectorial example with the same strict inequality.

More decisively, Theorem 6.4.4 of arXiv:1912.09347 states that for every
\(p\in(1,\infty)\setminus\{2\}\) and every \(\sigma\in(0,\pi)\), there is a
closed subspace \(Y\subset L^p([0,1])\) and a sectorial operator \(A\) on \(Y\)
such that \(A\) has a bounded \(H^\infty\)-calculus, has bounded imaginary
powers, is almost \(\gamma\)-sectorial, and
\[
  \omega(A)=0,
  \qquad
  \omega_{H^\infty}(A)=\sigma .
\]

\section*{Identification}

The notation \(\omega_H(A)\) in arXiv:math/0010155 is the same bounded
\(H^\infty\)-calculus angle denoted \(\omega_{H^\infty}(A)\) in the supporting
paper.  Since \(\sigma>0\) is arbitrary, Theorem 6.4.4 gives explicit
\(H^\infty\)-sectorial operators for which the two angles are unequal.
Therefore the source question has a negative literature answer.

\section*{Provenance and Scope}

This is a literature-implied answer rather than an original proof.  The
supporting paper does not explicitly frame Theorem 6.4.4 as answering the
Kalton--Weis sentence from arXiv:math/0010155; the implication follows by
matching the angle notation and hypotheses.  No broader claim is made here
about all remaining functional-calculus or maximal-regularity questions.

\section*{References}

\begin{itemize}
\item N. J. Kalton and L. Weis, \emph{The \(H^\infty\)-calculus and sums of
closed operators}, arXiv:math/0010155.
\item N. J. Kalton, E. Lorist, and L. Weis, \emph{Euclidean structures and
operator theory in Banach spaces}, arXiv:1912.09347.
\item N. J. Kalton, \emph{A remark on sectorial operators with an
\(H^\infty\)-calculus}, Contemp. Math. 321, 91--99, 2003.
\end{itemize}

\end{document}
